% ============================================================
%  ch10 —— §6（原文 p.40–p.44）
% ============================================================
\routenote{p.40--p.44}

\section{整体单值性}

本节致力于证明定理 3。设 $G$ 是 $k$ 上分裂几乎单群。

\subsection{对加性特征的依赖}

回忆注 2.7(3)：$T\rtimes\operatorname{Aut}^\dagger(G)\times\mathbb{G}_m^{\mathrm{rot}}$ 作用在 $\Bun_{G(0,2)}$ 与 Hecke 对应 (2.6) 上。群 $T(k)\rtimes\operatorname{Aut}^\dagger(G)\times\mathbb{G}_m^{\mathrm{rot}}(k)$ 也作用在 $I^\infty(1)/I^\infty(2)$ 上，从而作用在一般加性特征构成的空间上。记 $S_\varphi$ 为 $\varphi$ 在 $T\rtimes\operatorname{Aut}^\dagger(G)\times\mathbb{G}_m^{\mathrm{rot}}$ 作用下的稳定子。这是 $k$ 上的有限群概形。

当 $G$ 为伴随型时，对每个 $\sigma\in\operatorname{Aut}^\dagger(G)$，存在唯一的 $(t,s)\in(T\times\mathbb{G}_m^{\mathrm{rot}})(k)$ 使 $(t,\sigma,s)$ 固定 $\varphi$。因此在此情形，投影 $S_\varphi\to\operatorname{Aut}^\dagger(G)\xrightarrow{\sim}\operatorname{Out}(G)$ 是同构（作为 $k$ 上的离散群）。一般地，$S_\varphi\to\operatorname{Aut}^\dagger(G)$ 是 $Z_G$-挠子，但在 $k$-点上未必满。

下述引理直接由几何 Hecke 算子的定义推出。

\begin{lemma}
定义 $\Kl_{\ckG}(\varphi)$ 的张量函子
\[
\Hk_\varphi:\mathcal S\ni\mathrm{IC}_V\mapsto\Hk_V(A_\varphi)|_{\star\times\PP^1\setminus\{0,\infty\}}\in\Loc(\PP^1\setminus\{0,\infty\})
\]
带自然的 $S_\varphi$-等变结构。这里 $S_\varphi(k)$ 经其在 $\operatorname{Aut}^\dagger(G)$ 中的像作用在 $\mathcal S$ 上，并经 $S_\varphi\to\mathbb{G}_m^{\mathrm{rot}}$ 对 $\PP^1\setminus\{0,\infty\}$ 的作用作用在 $\Loc(\PP^1\setminus\{0,\infty\})$ 上。
\end{lemma}

在引理 9.3(1) 中我们将验证：在等价 $\Rep(\ckG)\cong\mathcal S$ 下，$\operatorname{Aut}^\dagger(G)$ 在 $\mathcal S$ 上的作用与 $\operatorname{Aut}^\dagger(G)\cong\operatorname{Aut}^\dagger(\ckG)$（相对于对偶群上选定的钉住）的作用一致。

记 $S^1_\varphi(k)=\ker(S_\varphi(k)\to\mathbb{G}_m^{\mathrm{rot}}(k))$，并记 $S^{\mathrm{rot}}_\varphi(k)$ 为 $S_\varphi(k)$ 在 $\mathbb{G}_m^{\mathrm{rot}}(k)$ 中的像，它是有限循环群。我们希望用上述等变性推出 $\Kl_{\ckG}(\varphi)$ 的单值表示可取值于 $\ckG^{S^1_\varphi(k)}$。然而，由于该表示只在内部自同构意义下有定义，这需要额外的论证。

回忆我们已选定 $0$ 处形式去心圆盘 $\operatorname{Spec}K_0$ 上方的几何一般点 $\eta$。我们已把 $\Kl_{\ckG}(\varphi,1)$ 定义为从 Satake 范畴 $\mathcal S\cong\Rep(\ckG)$ 到 $\Loc(\PP^1\setminus\{0,\infty\})$ 的函子，故限制到 $\eta$ 给出张量函子
\begin{equation}
\omega_\varphi:\mathcal S\xrightarrow{\ \Kl_{\ckG}(\varphi)\ }\Loc(\PP^1\setminus\{0,\infty\})\xrightarrow{j_\eta^*}\Vec.
\end{equation}
由于这是 $\mathcal S$ 的纤维函子，可定义 ${}_\varphi\ckG=\operatorname{Aut}^\otimes(\omega_\varphi)$。当然 ${}_\varphi\ckG$ 同构于 $\ckG$，但并非典范同构。群 ${}_\varphi\ckG$ 的好处是：由 Tannaka 形式化，得到典范同态 $\operatorname{Aut}^\otimes(j_\eta^*)\to\operatorname{Aut}^\otimes(\omega_\varphi)$，其中 $\operatorname{Aut}^\otimes(j_\eta^*)$ 是基本群 $\pi_1(\PP^1\setminus\{0,\infty\},\eta)$ 的 pro-代数包络。因此得到同态
\begin{equation}
\phi:\pi_1(\PP^1\setminus\{0,\infty\},\eta)\to{}_\varphi\ckG(\mathbb{Q}_\ell).
\end{equation}
由引理 6.1 得到同态
\begin{equation}
S^1_\varphi(k)\to\operatorname{Aut}^\otimes(\mathcal S,\Hk_\varphi)\xrightarrow{(j_\eta^*)_*}\operatorname{Aut}^\otimes(\mathcal S,\omega_\varphi)=\operatorname{Aut}({}_\varphi\ckG).
\end{equation}
换言之，$S^1_\varphi$ 作用在 ${}_\varphi\ckG$ 上。我们将在引理 9.4 中证明 ${}_\varphi\ckG$ 还带自然的钉住 $\ddagger$，且上述作用保持该钉住。钉住 $\ddagger$ 与 $\dagger$ 定义典范同构 $\mathrm{can}:{}_\varphi\ckG\xrightarrow{\sim}\ckG$。利用此等同，得到：

\begin{corollary}
单值表示 $\phi$ 扩充为正合列之间的同态：
\begin{equation}
\begin{tikzcd}[column sep=large, row sep=large]
\pi_1(\PP^1\setminus\{0,\infty\},\eta) \arrow[r, "{\#S^{\mathrm{rot}}_\varphi(k)\text{-}\mathrm{th}}"] \arrow[d, "\phi"'] & \pi_1(\PP^1\setminus\{0,\infty\},\eta) \arrow[r] \arrow[d, "\tilde\phi"] & S^{\mathrm{rot}}_\varphi(k) \arrow[d, equal] \\
\ckG^{S^1_\varphi(k)} \arrow[r, hook] & \ckG^{S^1_\varphi(k)}\rtimes S^{\mathrm{rot}}_\varphi(k) \arrow[r, two heads] & S^{\mathrm{rot}}_\varphi(k)
\end{tikzcd}
\end{equation}
\end{corollary}

\begin{proof}
由 (6.4)，$S^1_\varphi(k)$ 在 ${}_\varphi\ckG$ 上的作用经 $S^1_\varphi(k)\to\operatorname{Aut}^\otimes(\mathcal S,\Hk_\varphi)$ 分解。由引理 9.2，单值表示 $\phi$ 因而经 $\pi_1(\PP^1\setminus\{0,\infty\},\eta)\to{}_\varphi\ckG^{S^1_\varphi}\subset{}_\varphi\ckG$ 分解，即
\[
\Hk_\varphi:\mathcal S\cong\Rep({}_\varphi\ckG)\xrightarrow{\mathrm{Res}}\Rep({}_\varphi\ckG^{S^1_\varphi(k)})\xrightarrow{\kappa}\Loc(\PP^1\setminus\{0,\infty\})
\]
对某个张量函子 $\kappa$ 成立。

我们用钉住同构“$\mathrm{can}$”等同 $\ckG\cong{}_\varphi\ckG$，故可把函子 $\kappa$ 视为 $\Rep(\ckG^{S^1_\varphi(k)})\to\Loc(\PP^1\setminus\{0,\infty\})$。于是 $S^{\mathrm{rot}}_\varphi(k)=S_\varphi(k)/S^1_\varphi(k)$ 经 $S_\varphi(k)\to\operatorname{Out}(G)\xrightarrow{\iota}\operatorname{Aut}^\dagger(\ckG)$ 作用在 $\ckG^{S^1_\varphi(k)}$ 上，从而作用在 $\kappa$ 的源与靶上。引理 6.1 蕴涵 $\kappa$ 带自然的 $S^{\mathrm{rot}}_\varphi(k)$-等变结构。对两个张量范畴取 $S^{\mathrm{rot}}_\varphi(k)$-不变量，得到函子 $\tilde\kappa$：
\[
\begin{tikzcd}[column sep=large, row sep=large]
\Rep(\ckG^{S^1_\varphi(k)}\rtimes S^{\mathrm{rot}}_\varphi(k)) \arrow[r, "\tilde\kappa"] \arrow[d, "\wr"'] & \Loc(\PP^1\setminus\{0,\infty\}/S^{\mathrm{rot}}_\varphi(k)) \arrow[d, "\wr"] \\
\Rep(\ckG^{S^1_\varphi(k)})^{S^{\mathrm{rot}}_\varphi(k)} \arrow[r, "(\kappa)^{S^{\mathrm{rot}}_\varphi}"] & \Loc(\PP^1\setminus\{0,\infty\})^{S^{\mathrm{rot}}_\varphi(k)}
\end{tikzcd}
\]
由于商映射 $\PP^1\setminus\{0,\infty\}\to\PP^1\setminus\{0,\infty\}/S^{\mathrm{rot}}_\varphi(k)$ 可等同于 $\PP^1\setminus\{0,\infty\}=\mathbb{G}_m$ 的 $\#S^{\mathrm{rot}}_\varphi(k)$ 次幂映射，即得图 (6.6)。
\end{proof}

\begin{remark}
把 $k$ 换成扩域 $k'$、把 $\PP^1\setminus\{0,\infty\}$ 换成 $\PP^1\setminus\{0,\infty\}\otimes_kk'$ 后，推论 6.5 仍成立。特别地，为获得几何单值性的信息，取 $k'=\bar k$。
\end{remark}

当 $\operatorname{Out}(\ckG)$ 非平凡时，有

\medskip
\begin{table}[htbp]
\centering
\caption{外自同构与 $\varphi$ 的稳定子}
\label{tab:2}
\small
\begin{tabular}{ccccc}
\toprule
$\ckG$ & $\operatorname{Aut}^\dagger(\ckG)\xrightarrow{\sim}\operatorname{Out}(\ckG)$ & $\ckG^{\operatorname{Aut}^\dagger(\ckG)}$ & $S^1_\varphi(k)$ & $S^{\mathrm{rot}}_\varphi(k)$\\
\midrule
$A_{2n-1}$（$n\ge 2$） & $\mathbb{Z}/2$ & $C_n$ & $\operatorname{Out}(\ckG)$ & $1$\\
$A_{2n}$ & $\mathbb{Z}/2$ & $B_n$ & $1$ & $\mathbb{Z}/2$\\
$D_4$ & $S_3$ & $G_2$ & $\operatorname{Out}(\ckG)$ & $1$\\
$D_n$（$n\ge 5$） & $\mathbb{Z}/2$ & $B_{n-1}$ & $\operatorname{Out}(\ckG)$ & $1$\\
$E_6$ & $\mathbb{Z}/2$ & $F_4$ & $\operatorname{Out}(\ckG)$ & $1$\\
\bottomrule
\end{tabular}
\end{table}
\medskip

\subsection{整体单值性的 Zariski 闭包}

设 $\ckG^{\mathrm{geo}}\subset\ckG$ 是几何单值表示 $\phi^{\mathrm{geo}}:\pi_1(\PP^1\setminus\{0,\infty\}\otimes_k\bar k,\eta)\to{}_\varphi\ckG\cong\ckG$ 的像的 Zariski 闭包。先证明 $\ckG^{\mathrm{geo}}$ 不太小。

\begin{proposition}[B.~Gross]
若 $\operatorname{char}(k)>2$ 且 $G$ 的秩至少为 $2$，则 $\ckG^{\mathrm{geo}}$ 不包含于 $\ckG$ 的任何主 $\mathrm{PGL}_2$ 中。
\end{proposition}

\begin{proof}
反设 $\ckG^{\mathrm{geo}}\subset\mathrm{PGL}_2\subset\ckG$，其中 $\mathrm{PGL}_2$ 含 $\ckG$ 的一个主幂幺元（$I^t_0$ 的像）。野惯性群 $I^+_\infty$ 的像必须非平凡，因为 $\Kl^{\theta^\vee}_{\ckG}(\varphi)$ 在 $\infty$ 处有非零 Swan 导子。

由于 $p=\operatorname{char}(k)>2$，$\phi(I^+_\infty)$ 落在极大环面 $\mathbb{G}_m\subset\mathrm{PGL}_2$ 中，且含 $\mu_p\subset\mathbb{G}_m$。由于 $\phi(I_\infty)$ 正规化 $\phi(I^+_\infty)$，它必包含在环面 $\mathbb{G}_m$ 的正规化子 $N(\mathbb{G}_m)\subset\mathrm{PGL}_2$ 中。

对 $\mathrm{PGL}_2$ 的任一不可约表示 $S^{2\ell}=\operatorname{Sym}^{2\ell}(\mathrm{Std})$（$\mathrm{Std}$ 是 $\mathrm{SL}_2$ 的二维表示），每对权空间 $S^{2\ell}(n)\oplus S^{2\ell}(-n)$（$0\le n\le\ell$）在 $N(\mathbb{G}_m)$ 下稳定，从而在 $I_\infty$ 下稳定。若权 $n$ 不整除 $p$，则 $S^{2\ell}(n)\oplus S^{2\ell}(-n)$ 的 Swan 导子至少为 $1$。因此
\[
\Swan_{I_\infty}(S^{2\ell})\ge\ell-[\ell/p].
\]
考虑 $I_\infty$ 在 $\ckG$ 的拟极小表示 $V_{\theta^\vee}$ 上的作用。由引理 10.1，$V_{\theta^\vee}$ 分解为主 $\mathrm{PGL}_2$ 的 $r_s(\ckG)$ 个不可约表示：$V_{\theta^\vee}=S^{2\ell_1}\oplus\cdots\oplus S^{2\ell_{r_s(\ckG)}}$。因此
\[
\Swan_{I_\infty}(V_{\theta^\vee})\ge\sum_{i=1}^{r_s(\ckG)}(\ell_i-[\ell_i/p])\ge\Big(1-\frac1p\Big)\frac{h\,r_s(\ckG)}{2}.
\]
这里用了 $\sum_i\ell_i=hr_s(\ckG)/2$（见引理 10.1 的证明）。故只要 $h>3$ 或 $p>3$，即得矛盾。即使在 $h=3$、$p=3$ 时，此时 $\ckG=\mathrm{PGL}_3$、$V_{\theta^\vee}=S^4\oplus S^2$，仍有 $\Swan(V_{\theta^\vee})\ge 2+1=3>r_s(\ckG)$，矛盾！
\end{proof}

\begin{proof}[定理 3 的证明]
由于对每个 $V\in\Rep(\ckG)$，$\Kl^V_{\ckG}(\varphi)$ 是权 $0$ 的纯层，故几何半单。由 Deligne [11, Corollaire 1.3.9]，$\ckG^{\mathrm{geo}}$ 的中性分支 $\ckG^\circ_{\mathrm{geo}}$ 是半单群。

\emph{第一步.} 设 $G$ 型不为 $A_1$、$A_{2n}$ 或 $B_3$。先确定 $\ckG^\circ_{\mathrm{geo}}$，等价地确定其 Lie 代数 $\ckg_{\mathrm{geo}}$。

一方面，由定理 1(2)，$\ckG^\circ_{\mathrm{geo}}$ 含主幂幺元，故含一个主 $\mathrm{PGL}_2$。命题 6.8 说 $\ckG^{\mathrm{geo}}$ 不能等于 $\mathrm{PGL}_2$。由于主 $\mathrm{PGL}_2$ 在 $\ckG$ 中是自正规化的，故 $\ckG^\circ_{\mathrm{geo}}$ 不能等于 $\mathrm{PGL}_2$，即 $\mathfrak{sl}_2\subsetneq\ckg_{\mathrm{geo}}$。

另一方面，由于 $G$ 型不为 $A_{2n}$，$S^1_\varphi(k)=\operatorname{Aut}^\dagger(\ckG)\xrightarrow{\sim}\operatorname{Out}(\ckG)$。在此情形，推论 6.5 蕴涵 $\ckG^{\mathrm{geo}}\subset\ckG^{\operatorname{Aut}^\dagger(\ckG)}$。

Dynkin 分类了 $\ckg$ 中含主 $\mathfrak{sl}_2$ 的所有 Lie 子代数。若 $\ckg\neq\mathfrak{so}_7$，则或者 $\ckg_{\mathrm{geo}}=\mathfrak{sl}_2$，或者 $\ckg_{\mathrm{geo}}$ 是 $\ckG$ 的某个钉住自同构的不动点代数。在我们的情形，若 $G$ 型不为 $A_{2n}$，已可断定 $\ckg_{\mathrm{geo}}=\ckg^{\operatorname{Aut}^\dagger(\ckG)}$。

\emph{第二步.} 设 $\ckG$ 型为 $B_3$ 且 $\operatorname{char}(k)>3$，断言 $\ckG^\circ_{\mathrm{geo}}=G_2$。为此只需证明 $\Kl^{\wedge^3V_7}_{\mathrm{SO}_7}(\varphi)$ 在 $\PP^1\setminus\{0,\infty\}\otimes_k\bar k$ 上有整体截面，其中 $V_7$ 是 $\mathrm{SO}_7$ 的 $7$ 维标准表示。设不然，则把 (5.4)–(5.6) 的长正合列中的 $\ckg$ 换成 $\wedge^3V_7$，将推出
\begin{equation}
\dim(\wedge^3V_7)^{I_0}+\dim(\wedge^3V_7)^{I_\infty}\le\dim H^1_c(\PP^1\setminus\{0,\infty\},\Kl^{\wedge^3V_7}_{\mathrm{SO}_7}(\varphi))=\Swan_{I_\infty}(\wedge^3V_7).
\end{equation}
对 $V_7$ 取标准基 $\{e_{-3},\dots,e_0,\dots,e_3\}$（二次型为 $e_0^2+e_1e_{-1}+e_2e_{-2}+e_3e_{-3}$）。$\wedge^3V_7$ 在主 $\mathfrak{sl}_2$ 下的 $0$-权空间由 $\{e_1\wedge e_{-2}\wedge e_{-3},\ e_{-1}\wedge e_2\wedge e_3,\ e_0\wedge e_i\wedge e_{-i}\ (i=1,2,3)\}$ 张成，故维数为 $5$。由定理 1(3)，$I_0$ 作为主幂幺元作用在 $V_7$ 上，故 $\dim(\wedge^3V_7)^{I_0}=5$。

由于 $\operatorname{char}(k)>3$ 不整除 $\ckG$ 的 Coxeter 数，推论 2.13 可用。由于 $I^+_\infty$ 在 $V_7$ 上作用的 breaks 为 $1/h=1/6$，$I^+_\infty$ 在 $\wedge^3V_7$ 的各非零权空间上的 breaks 至多为 $1/6$。因此 $\Swan_{I_\infty}(\wedge^3V_7)\le[\,\tfrac16\dim\wedge^3V_7\,]=[35/6]=5$。此外，由推论 2.13 中 $\phi(I^+_\infty)$ 的描述，$I^+_\infty$ 平凡地作用在 $\operatorname{Span}\{e_1\wedge e_2\wedge e_3,\ e_{-1}\wedge e_{-2}\wedge e_{-3}\}$ 上，而 Coxeter 置换 $e_1\to e_2\to e_3\to e_{-1}\to e_{-2}\to e_{-3}\to e_1$ 置换 $e_1\wedge e_2\wedge e_3$ 与 $e_{-1}\wedge e_{-2}\wedge e_{-3}$。故 $\dim(\wedge^3V_7)^{I_\infty}\ge 1$。于是
\[
\dim(\wedge^3V_7)^{I_0}+\dim(\wedge^3V_7)^{I_\infty}\ge 5+1>5\ge\Swan_{I_\infty}(\wedge^3V_7),
\]
与 (6.9) 矛盾。这证明当 $\ckG=\mathrm{SO}_7$ 时 $\ckG^{\mathrm{geo}}=G_2$。

\emph{第三步.} 当 $\ckG$ 型为 $A_{2n}$ 时，$\Kl_{\ckG}(\varphi)$ 来自经典 Kloosterman 层 $\mathrm{Kl}_n(\varphi)$，其整体单值性由 Katz 在 [25] 中处理，见 (0.2)。

\emph{第四步.} 剩下证明 $\ckG^{\mathrm{geo}}$ 连通。$A_n$ 的情形由 Katz 在上引文献中处理。$\ckG$ 型为 $E_6$ 时，容易验证 $\ckG^{\operatorname{Aut}^\dagger(\ckG)}$ 已是连通的。

一般地有 $\ckG^{\mathrm{geo}}\subset N_{\ckG}(\ckG^\circ_{\mathrm{geo}})=\ckG^\circ_{\mathrm{geo}}Z_{\ckG}$。考虑满同态
\[
\phi^{\mathrm{geo}}:\pi_1^{\mathrm{geo}}(\PP^1\setminus\{0,\infty\},\eta)\xrightarrow{\phi^{\mathrm{geo}}}\ckG^{\mathrm{geo}}\to\pi_0(\ckG^{\mathrm{geo}})=Z_{\ckG}/Z_{\ckG}\cap\ckG^\circ_{\mathrm{geo}}.
\]
假设 $\ckG$ 型不为 $A_n$ 或 $E_6$，则 $\pi_0(\ckG)$ 是 $2$-群。由于 $\operatorname{char}(k)>2$，$\phi^{\mathrm{geo}}$ 经驯商分解。另一方面，$I_0\cong\pi_1^{\mathrm{geo}}(\PP^1\setminus\{0,\infty\},\eta)$ 中的驯生成元必须映到 $\ckG^{\mathrm{geo}}$ 中的幂幺元，从而落在 $\ckG^\circ_{\mathrm{geo}}$ 内，因此映射 $\phi^{\mathrm{geo}}$ 平凡，即 $\ckG^{\mathrm{geo}}$ 连通。
\end{proof}
